\documentclass[a4paper,11pt]{article}

\usepackage[utf8]{inputenc}
\usepackage[T1]{fontenc}
\usepackage[ngerman]{babel}
\usepackage{amsmath,amssymb}
\usepackage{geometry}
\geometry{margin=2.0cm}
\usepackage{enumitem}
\usepackage{fancyhdr}
\usepackage{xcolor}
\usepackage{tikz}
\usepackage{titlesec}
\usepackage{array}
\usepackage{colortbl}

%================= Dark-Mode Farben (Catppuccin Mocha) =================
\definecolor{base}{HTML}{1E1E2E}
\definecolor{fg}{HTML}{CDD6F4}
\definecolor{accent}{HTML}{89B4FA}
\definecolor{accent2}{HTML}{F5C2E7}
\definecolor{gridcolor}{HTML}{45475A}

\pagecolor{base}

\newcommand{\transp}{^{\mathsf{T}}}
\newcommand{\norm}[1]{\left\lVert #1 \right\rVert}
\newcommand{\inner}[2]{\left\langle #1,\, #2 \right\rangle}
\newcommand{\rang}{\operatorname{rang}}
\newcommand{\spur}{\operatorname{spur}}
\newcommand{\diag}{\operatorname{diag}}
\newenvironment{psmallmatrix}{\left(\begin{smallmatrix}}{\end{smallmatrix}\right)}

\newcommand{\pkt}[1]{\hfill\textnormal{\color{accent2}\itshape (#1)}}

\newcommand{\karo}[1]{\par\vspace{0.3cm}\noindent
  \begin{tikzpicture}
    \draw[step=5mm, gridcolor, very thin] (0,0) grid (\linewidth,#1);
  \end{tikzpicture}\par\vspace{0.3cm}}

\newcommand{\kasten}{{\color{gridcolor}$\square$}}

\titleformat{\section}{\color{accent}\normalfont\Large\bfseries}{}{0pt}{}

\pagestyle{fancy}
\fancyhf{}
\lhead{\color{fg}\small FLA \& Analytische Geometrie}
\rhead{\color{fg}\small Probeklausur 02 (Vollumfang)}
\cfoot{\color{fg}\thepage}
\renewcommand{\headrule}{{\color{gridcolor}\hrule height 0.4pt}}

\arrayrulecolor{gridcolor}
\setlength{\parindent}{0pt}

\begin{document}
\color{fg}

\thispagestyle{fancy}

\begin{center}
  {\color{accent}\LARGE\bfseries Probeklausur (Vollumfang)}\\[0.3em]
  {\color{fg}\large Fortgeschrittene Lineare Algebra und Analytische Geometrie}\\[0.6em]
  {\color{gridcolor}\rule{\linewidth}{0.4pt}}
\end{center}

\vspace{0.4em}

\renewcommand{\arraystretch}{1.6}
\begin{tabular}{@{}p{0.5\linewidth}@{}p{0.5\linewidth}@{}}
  Name: \ \rule[-2pt]{0.7\linewidth}{0.4pt} &
  Datum: \ \rule[-2pt]{0.55\linewidth}{0.4pt}
\end{tabular}

\vspace{0.6em}

\begin{tabular}{@{}l p{0.72\linewidth}@{}}
  \textbf{\color{accent}Bearbeitungszeit:} & 150 Minuten \\
  \textbf{\color{accent}Gesamtpunkte:}     & 130 Punkte \\
  \textbf{\color{accent}Hilfsmittel:}      & beidseitig handbeschriebenes Formelblatt (DIN A4), nicht-programmierbarer Taschenrechner \\
\end{tabular}

\vspace{0.4em}

\textit{\color{fg}\small Hinweis: Diese Klausur deckt den \textbf{gesamten} Stoff der Open-Book-Formelsammlung ab (alle Zerlegungen, Eigenwerte/SVD, Normen \& innere Produkte, Geometrie/Hyperebenen, PCA, Regression, Finite Differenzen). Alle Ergebnisse sind zu begründen.}

\vspace{0.9em}

\renewcommand{\arraystretch}{1.5}
\noindent
\begin{tabular}{|c|*{15}{c|}c|}
\hline
\scriptsize\textbf{Aufg.} & 1 & 2 & 3 & 4 & 5 & 6 & 7 & 8 & 9 & 10 & 11 & 12 & 13 & 14 & 15 & \textbf{$\Sigma$} \\
\hline
\scriptsize\textbf{max.} & 8 & 8 & 10 & 9 & 8 & 9 & 9 & 12 & 8 & 8 & 8 & 8 & 8 & 8 & 9 & \textbf{130} \\
\hline
\scriptsize\textbf{erm.} & & & & & & & & & & & & & & & & \\
\hline
\end{tabular}

\renewcommand{\arraystretch}{1.0}

% ================================================================
\clearpage
\section*{Aufgabe 1 -- Quickies \pkt{8 Punkte}}
Sind die folgenden Aussagen \emph{wahr} oder \emph{falsch}? Kreuzen Sie an.
\vspace{0.5em}
\begin{enumerate}[label=\alph*)]
  \item Jede reelle symmetrische Matrix ist orthogonal diagonalisierbar. \hfill \kasten\ wahr \quad \kasten\ falsch
  \item Für jede Matrix $A$ ist $A\transp A$ positiv semidefinit. \hfill \kasten\ wahr \quad \kasten\ falsch
  \item Die Spur einer Matrix ist gleich der Summe ihrer Eigenwerte. \hfill \kasten\ wahr \quad \kasten\ falsch
  \item Eine nicht durch den Ursprung verlaufende Ebene ist ein Untervektorraum. \hfill \kasten\ wahr \quad \kasten\ falsch
  \item Für die Maximumsnorm gilt $\norm{x}_\infty = \max_i |x_i|$. \hfill \kasten\ wahr \quad \kasten\ falsch
  \item Die Singulärwerte von $A$ sind die Eigenwerte von $A$. \hfill \kasten\ wahr \quad \kasten\ falsch
  \item Für invertierbare $A,B$ gilt $(AB)^{-1} = A^{-1}B^{-1}$. \hfill \kasten\ wahr \quad \kasten\ falsch
  \item Bei der CR-Zerlegung gilt $\dim C(A) = \dim C(A\transp)$. \hfill \kasten\ wahr \quad \kasten\ falsch
\end{enumerate}
\karo{6cm}

% ================================================================
\clearpage
\section*{Aufgabe 2 -- Lineares Gleichungssystem (Gauß) \pkt{8 Punkte}}
Gegeben sei das Gleichungssystem
\[
\begin{aligned}
 x_1 + \phantom{2}x_2 + 2x_3 &= 4,\\
2x_1 + 3x_2 + 3x_3 &= 9,\\
3x_1 + 4x_2 + 5x_3 &= 13.
\end{aligned}
\]
\begin{enumerate}[label=\alph*)]
  \item Bringen Sie das System mit dem Gauß-Verfahren auf Zeilenstufenform. \pkt{4}
  \item Bestimmen Sie die \textbf{allgemeine} Lösung (Partikulärlösung $+$ Nullraumanteil). \pkt{4}
\end{enumerate}
\karo{13cm}

% ================================================================
\clearpage
\section*{Aufgabe 3 -- CR-Zerlegung \& vier Unterräume \pkt{10 Punkte}}
Gegeben sei
\[
A = \begin{pmatrix} 1 & 2 & 0 & 1 \\ 2 & 4 & 1 & 4 \\ 1 & 2 & 1 & 3 \end{pmatrix}.
\]
\begin{enumerate}[label=\alph*)]
  \item Bestimmen Sie $\rang(A)$ und die CR-Zerlegung $A = CR$. \pkt{4}
  \item Geben Sie eine Basis des Spaltenraums $C(A)$ an. \pkt{2}
  \item Bestimmen Sie eine Basis des Nullraums $N(A)$. \pkt{2}
  \item Geben Sie $\dim C(A),\ \dim C(A\transp),\ \dim N(A),\ \dim N(A\transp)$ an. \pkt{2}
\end{enumerate}
\karo{14cm}

% ================================================================
\clearpage
\section*{Aufgabe 4 -- LU-Zerlegung \& LGS lösen \pkt{9 Punkte}}
Gegeben seien
\[
A = \begin{pmatrix} 1 & 2 & 1 \\ 2 & 5 & 3 \\ 3 & 8 & 9 \end{pmatrix},
\qquad
b = \begin{pmatrix} 4 \\ 10 \\ 20 \end{pmatrix}.
\]
\begin{enumerate}[label=\alph*)]
  \item Bestimmen Sie eine LU-Zerlegung $A = LU$. \pkt{4}
  \item Bestimmen Sie $\det(A)$. \pkt{1}
  \item Lösen Sie $Ax = b$ mittels Vorwärts- ($Ly=b$) und Rückwärtssubstitution ($Ux=y$). \pkt{4}
\end{enumerate}
\karo{13cm}

% ================================================================
\clearpage
\section*{Aufgabe 5 -- Cholesky-Zerlegung \& Lösen \pkt{8 Punkte}}
Gegeben seien
\[
A = \begin{pmatrix} 4 & 2 & 4 \\ 2 & 10 & 5 \\ 4 & 5 & 6 \end{pmatrix},
\qquad
b = \begin{pmatrix} 8 \\ 7 \\ 10 \end{pmatrix}.
\]
\begin{enumerate}[label=\alph*)]
  \item Bestimmen Sie die Cholesky-Zerlegung $A = LL\transp$ (und begründen Sie damit die positive Definitheit). \pkt{5}
  \item Lösen Sie $Ax = b$ über $Ly = b$ und $L\transp x = y$. \pkt{3}
\end{enumerate}
\karo{14cm}

% ================================================================
\clearpage
\section*{Aufgabe 6 -- QR-Zerlegung (Gram-Schmidt) \pkt{9 Punkte}}
Gegeben seien
\[
A = \begin{pmatrix} 3 & 4 \\ 4 & 3 \end{pmatrix},
\qquad
b = \begin{pmatrix} 7 \\ 7 \end{pmatrix}.
\]
\begin{enumerate}[label=\alph*)]
  \item Bestimmen Sie mit dem Gram-Schmidt-Verfahren die QR-Zerlegung $A = QR$. \pkt{6}
  \item Lösen Sie $Ax = b$ über $Rx = Q\transp b$. \pkt{3}
\end{enumerate}
\karo{14cm}

% ================================================================
\clearpage
\section*{Aufgabe 7 -- Eigenwerte, Diagonalisierung \& Matrixpotenz \pkt{9 Punkte}}
Gegeben sei
\[
A = \begin{pmatrix} 3 & 1 \\ 1 & 3 \end{pmatrix}.
\]
\begin{enumerate}[label=\alph*)]
  \item Bestimmen Sie die Eigenwerte. \pkt{3}
  \item Bestimmen Sie Eigenvektoren sowie $V$ und $D$ mit $A = V D V^{-1}$. \pkt{3}
  \item Berechnen Sie $A^2$ über die Diagonalisierung $A^2 = V D^2 V^{-1}$. \pkt{3}
\end{enumerate}
\karo{13cm}

% ================================================================
\clearpage
\section*{Aufgabe 8 -- SVD, Pseudoinverse \& Rang-1-Näherung \pkt{12 Punkte}}
Gegeben sei
\[
A = \begin{pmatrix} 1 & 2 \\ 2 & 1 \end{pmatrix}.
\]
\begin{enumerate}[label=\alph*)]
  \item Berechnen Sie $A\transp A$ und dessen Eigenwerte. \pkt{2}
  \item Bestimmen Sie die Singulärwerte $\sigma_1 \ge \sigma_2$. \pkt{2}
  \item Bestimmen Sie $V$ und $U$, sodass $A = U\Sigma V\transp$. \pkt{4}
  \item Bestimmen Sie die Pseudoinverse $A^{+}$. \pkt{2}
  \item Bestimmen Sie die beste Rang-1-Näherung $A_1 = \sigma_1 u_1 v_1\transp$. \pkt{2}
\end{enumerate}
\karo{14cm}

% ================================================================
\clearpage
\section*{Aufgabe 9 -- Determinante, Inverse, Spur \& Rechenregeln \pkt{8 Punkte}}
\begin{enumerate}[label=\alph*)]
  \item Berechnen Sie $\det(M)$ für $M = \begin{pmatrix} 2 & 1 & 0 \\ 1 & 3 & 1 \\ 0 & 1 & 2 \end{pmatrix}$ (z.\,B.\ mit Sarrus). \pkt{3}
  \item Berechnen Sie die Inverse von $N = \begin{pmatrix} 4 & 3 \\ 5 & 4 \end{pmatrix}$. \pkt{2}
  \item Bestimmen Sie $\spur(M)$ und erläutern Sie den Zusammenhang zu den Eigenwerten. \pkt{1}
  \item Zeigen Sie: Für jede Matrix $B$ ist $B\transp B$ symmetrisch. \pkt{2}
\end{enumerate}
\karo{13cm}

% ================================================================
\clearpage
\section*{Aufgabe 10 -- Normen \& innere Produkte \pkt{8 Punkte}}
Gegeben seien
\[
x = \begin{pmatrix} 2 \\ -1 \\ 2 \end{pmatrix},
\qquad
y = \begin{pmatrix} 1 \\ 2 \\ 2 \end{pmatrix}.
\]
\begin{enumerate}[label=\alph*)]
  \item Berechnen Sie $\norm{x}_1$, $\norm{x}_2$ und $\norm{x}_\infty$. \pkt{3}
  \item Berechnen Sie $\inner{x}{y}$ und den Winkel zwischen $x$ und $y$. \pkt{3}
  \item Überprüfen Sie die Cauchy-Schwarz-Ungleichung $|\inner{x}{y}| \le \norm{x}_2\,\norm{y}_2$. \pkt{2}
\end{enumerate}
\karo{12cm}

% ================================================================
\clearpage
\section*{Aufgabe 11 -- Orthogonale Matrizen \& Untervektorraum \pkt{8 Punkte}}
\begin{enumerate}[label=\alph*)]
  \item Prüfen Sie, ob
  $Q = \dfrac{1}{3}\begin{pmatrix} 2 & -1 & 2 \\ 2 & 2 & -1 \\ -1 & 2 & 2 \end{pmatrix}$
  orthogonal ist. \pkt{4}
  \item Ist $U = \{\, x \in \mathbb{R}^3 : x_1 + x_2 + x_3 = 0 \,\}$ ein Untervektorraum von $\mathbb{R}^3$? Und $W = \{\, x \in \mathbb{R}^3 : x_1 + x_2 + x_3 = 1 \,\}$? Begründen Sie. \pkt{4}
\end{enumerate}
\karo{13cm}

% ================================================================
\clearpage
\section*{Aufgabe 12 -- Hyperebene \& Abstände \pkt{8 Punkte}}
Gegeben sei die Hyperebene
\[
2x_1 + x_2 - 2x_3 = 6
\]
sowie die Punkte $p = (3,1,1)\transp$ und $q = (1,2,3)\transp$.
\begin{enumerate}[label=\alph*)]
  \item Bestimmen Sie einen Normalenvektor und $\norm{n}$. \pkt{1}
  \item Liegt $p$ auf der Ebene? Auf welcher Seite liegt $p$? \pkt{2}
  \item Berechnen Sie den Abstand von $p$ zur Ebene. \pkt{2}
  \item Berechnen Sie den Abstand der Punkte $p$ und $q$. \pkt{3}
\end{enumerate}
\karo{12cm}

% ================================================================
\clearpage
\section*{Aufgabe 13 -- Lineare Regression \pkt{8 Punkte}}
Gegeben seien die Datenpunkte
\[
(1,1),\quad (2,2),\quad (3,2).
\]
Gesucht ist die Ausgleichsgerade $y = c_0 + c_1 x$.
\begin{enumerate}[label=\alph*)]
  \item Stellen Sie die Designmatrix $X$ und die Normalengleichung $X\transp X\,c = X\transp y$ auf. \pkt{4}
  \item Lösen Sie die Normalengleichung und geben Sie $c_0, c_1$ an. \pkt{4}
\end{enumerate}
\karo{13cm}

% ================================================================
\clearpage
\section*{Aufgabe 14 -- Hauptkomponentenanalyse (PCA) \pkt{8 Punkte}}
Gegeben seien die Datenpunkte
\[
(5,4),\quad (4,5),\quad (1,2),\quad (2,1).
\]
\begin{enumerate}[label=\alph*)]
  \item Bestimmen Sie den Mittelwert und die zentrierten Daten $A_c$. \pkt{3}
  \item Bestimmen Sie die Kovarianzmatrix $C = \tfrac{1}{n-1} A_c\transp A_c$. \pkt{3}
  \item Bestimmen Sie die erste Hauptkomponente (Richtung des größten Eigenwerts). \pkt{2}
\end{enumerate}
\karo{13cm}

% ================================================================
\clearpage
\section*{Aufgabe 15 -- Finite Differenzen \pkt{9 Punkte}}
Gegeben sei das Randwertproblem
\[
-u''(x) = 2, \qquad u(0) = u(1) = 0,
\]
mit den inneren Gitterpunkten $x_1 = 0{,}25,\ x_2 = 0{,}5,\ x_3 = 0{,}75$.
\begin{enumerate}[label=\alph*)]
  \item Leiten Sie mit dem zentralen Differenzenquotienten das lineare Gleichungssystem $T u = h^2 f$ her ($T$ und rechte Seite angeben). \pkt{5}
  \item Lösen Sie das Gleichungssystem (nutzen Sie die Symmetrie). \pkt{4}
\end{enumerate}
\karo{14cm}

\end{document}
